\chapter{Objetivos y preguntas de investigación}
\label{cap:objetivos}

\section{Objetivo general}

El objetivo general de este trabajo es determinar, con evidencia empírica y contraste estadístico,
si el aprendizaje federado constituye una alternativa útil para la previsión de demanda en
distribución alimentaria cuando la centralización de datos no es viable, y cuantificar el valor de
negocio que aportaría su adopción.

La formulación es deliberadamente neutra respecto al resultado. El trabajo no parte de la hipótesis
de que el aprendizaje federado sea superior, sino de la pregunta de si lo es, y admite como
resultado válido que no lo sea.

\section{Preguntas de investigación}

El objetivo general se concreta en cinco preguntas, dos técnicas, dos de negocio y una de alcance
ampliado.

\begin{description}
  \item[RQ1 (técnica).] ¿Un modelo federado predice mejor que el modelo local de cada tienda y se
    aproxima al centralizado, sin compartir datos crudos, en presencia de heterogeneidad entre
    silos? La pregunta incluye la personalización posterior del modelo federado como variante a
    evaluar.

  \item[RQ2 (técnica).] ¿Supera el modelo federado a los métodos convencionales de previsión en
    distribución, es decir, a los baselines ingenuos, al suavizado exponencial y a los árboles
    potenciados por gradiente?

  \item[RQ3 (negocio).] ¿Qué valor económico genera la mejora de precisión atribuible a la
    colaboración federada, expresado como reducción de error de previsión y traducido a un rango
    monetario?

  \item[RQ4 (negocio).] ¿Por qué sería preferible un esquema federado a un \emph{data clean room}
    para la previsión colaborativa entre operadores que no comparten datos?

  \item[RQ5 (ampliación).] ¿Qué coste en precisión tendría añadir garantías formales de privacidad
    mediante privacidad diferencial?
\end{description}

\section{Objetivos específicos}

Para responder a esas preguntas se definen los siguientes objetivos específicos:

\begin{enumerate}
  \item Construir un conjunto de datos de previsión a partir de un histórico real de ventas de
    distribución, con una partición temporal que impida la fuga de información y una definición de
    silos que reproduzca heterogeneidad realista entre participantes.

  \item Implementar un modelo de previsión común a todas las condiciones experimentales, de modo
    que las diferencias observadas sean atribuibles a la forma de entrenamiento y no a la
    arquitectura.

  \item Evaluar cinco condiciones —local, centralizado por silo, centralizado global, federado y
    federado con personalización— y compararlas entre sí y frente a los métodos convencionales
    mediante un protocolo único y un contraste estadístico pareado.

  \item Desplegar el sistema federado sobre infraestructura distribuida real, con aislamiento
    efectivo de los datos de cada silo, para verificar que el algoritmo estudiado en simulación se
    comporta igual fuera de ella.

  \item Cuantificar el valor económico de la mejora obtenida y compararlo con la alternativa de los
    \emph{data clean rooms}.
\end{enumerate}

\section{Alcance y delimitación}

El trabajo se delimita en los siguientes términos. La granularidad de previsión es tienda ×
familia de producto × semana, no referencia individual ni día, por las razones que se detallan en
el \cref{cap:materiales}. El horizonte de evaluación son las ocho últimas semanas del histórico
disponible. El modelo federado se entrena con los silos como participantes, no con las tiendas
individuales, lo que corresponde al escenario \emph{cross-silo}.

La pregunta RQ5, relativa al coste en precisión de la privacidad diferencial, se plantea como
ampliación y no se aborda experimentalmente en este trabajo; se retoma en las líneas de trabajo
futuro del \cref{cap:conclusiones}, apoyada en la recomendación regulatoria descrita en el
\cref{cap:estado_arte}.
